\lecture{7.4}{AIPTW Implementation in R}
\label{lec:7-4}

\subsection{Fitting the Outcome Regression}

To compute AIPTW, first estimate the outcome regression
\[
\Qbar(a,w)=\E(Y\given A=a,W=w).
\]
For a binary outcome, logistic regression is a simple starting point:

\begin{implementationbox}
\begin{lstlisting}[language=R]
# dat contains Y, A, and W variables
or_fit <- glm(Y ~ A + W, data = dat, family = binomial())

dat1 <- dat
dat1$A <- 1
Qbar_1W <- predict(or_fit, newdata = dat1, type = "response")

dat0 <- dat
dat0$A <- 0
Qbar_0W <- predict(or_fit, newdata = dat0, type = "response")
\end{lstlisting}
\end{implementationbox}

The objects \texttt{Qbar\_1W} and \texttt{Qbar\_0W} are estimates of
$\Qbar(1,W_i)$ and $\Qbar(0,W_i)$ for every observed individual.

\subsection{Fitting the Propensity Score}

Next estimate the propensity score:
\[
g_1(w)=\Prob(A=1\given W=w).
\]

\begin{implementationbox}
\begin{lstlisting}[language=R]
ps_fit <- glm(A ~ W, data = dat, family = binomial())

g_1W <- predict(ps_fit, newdata = dat, type = "response")
g_0W <- 1 - g_1W
\end{lstlisting}
\end{implementationbox}

In practice, both $\Qbar_n$ and $g_n$ can be estimated using Super Learner. The
subsequent AIPTW calculations are the same once predicted values are available.

\subsection{Computing AIPTW for $\E\{\Yone\}$}

For treatment level $1$,
\[
\psi_{n,\AIPTW}(1)
=
\frac{1}{n}\sum_{i=1}^n
\left[
\Qbar_n(1,W_i)
+
\frac{\ind(A_i=1)}{g_{n,1}(W_i)}
\{Y_i-\Qbar_n(1,W_i)\}
\right].
\]
In R:

\begin{implementationbox}
\begin{lstlisting}[language=R]
A <- dat$A
Y <- dat$Y

aiptw_ey1 <- mean(Qbar_1W + (A == 1) / g_1W * (Y - Qbar_1W))
\end{lstlisting}
\end{implementationbox}

This expression is often safer than writing separate means for the g-computation and
augmentation terms because it keeps the individual contribution visible.

\subsection{Computing AIPTW for $\E\{\Yzero\}$}

For treatment level $0$,
\[
\psi_{n,\AIPTW}(0)
=
\frac{1}{n}\sum_{i=1}^n
\left[
\Qbar_n(0,W_i)
+
\frac{\ind(A_i=0)}{g_{n,0}(W_i)}
\{Y_i-\Qbar_n(0,W_i)\}
\right].
\]
In R:

\begin{implementationbox}
\begin{lstlisting}[language=R]
aiptw_ey0 <- mean(Qbar_0W + (A == 0) / g_0W * (Y - Qbar_0W))
\end{lstlisting}
\end{implementationbox}

\subsection{Computing the ATE}

The AIPTW ATE estimator is
\[
\widehat{\ATE}_{\AIPTW}
=
\psi_{n,\AIPTW}(1)-\psi_{n,\AIPTW}(0).
\]

\begin{implementationbox}
\begin{lstlisting}[language=R]
aiptw_ate <- aiptw_ey1 - aiptw_ey0
\end{lstlisting}
\end{implementationbox}

Equivalently, define individual ATE contributions
\[
D_{i,n}
=
\Qbar_n(1,W_i)-\Qbar_n(0,W_i)
+
\frac{A_i}{g_{n,1}(W_i)}\{Y_i-\Qbar_n(1,W_i)\}
-
\frac{1-A_i}{g_{n,0}(W_i)}\{Y_i-\Qbar_n(0,W_i)\}.
\]
Then
\[
\widehat{\ATE}_{\AIPTW}=\frac{1}{n}\sum_{i=1}^nD_{i,n}.
\]

\subsection{Practical Coding Workflow}

\begin{estimationbox}
AIPTW coding workflow:
\begin{enumerate}
\item Fit or otherwise estimate $\Qbar_n(A,W)$.
\item Predict $\Qbar_n(1,W_i)$ and $\Qbar_n(0,W_i)$ for every individual.
\item Fit or otherwise estimate $g_{n,1}(W)$.
\item Compute $g_{n,0}(W_i)=1-g_{n,1}(W_i)$.
\item Form $\psi_{n,\AIPTW}(1)$ and $\psi_{n,\AIPTW}(0)$.
\item Subtract to estimate the ATE.
\item Compute influence-function-based standard errors as described in the next
lecture.
\end{enumerate}
\end{estimationbox}

Always check the range of estimated propensity scores. Extremely small values of
$g_{n,1}(W_i)$ among treated observations or $g_{n,0}(W_i)$ among untreated
observations can create unstable residual corrections.

\begin{implementationbox}
\begin{lstlisting}[language=R]
summary(g_1W)
summary(g_0W)
summary((A == 1) / g_1W)
summary((A == 0) / g_0W)
\end{lstlisting}
\end{implementationbox}
